% Claims and numbers for this section: Seshat-Star/docs/CLAIMS.md §4; RESULTS.md §1.4, §3.4, §4; narrative: docs/PAPER_FRAMING.md
% v1 drafted 11-16.  Number provenance (check before freeze):
% * Blind vs aware reader 16.5 -> 17.5 on originals, n = 200 pairs: RESULTS §1.4 (10-21 run, older verifier).
% * Wrong-target probe 2 of 15: RESULTS §3.4 (scene0549_00_ann39 is the inspected case). n is small; do not quote an interval.
% * Human check (tab:human, and 75 / 79 / 12 / 51 in the text): E7-A1 SNAPSHOT, analysis/_e7_snapshot.py.  12% / 51% = scorer
%   certifies 23 / 101 of the 200 A1 items (a1_items.jsonl, field `ours`).  79% = P(second annotator resolves | first resolves) = 146/186.
% * Predicate cards (E7-B SNAPSHOT, 218 judgements): three self-defined families 114/146 = 78%; stacks 53/61 = 87%; rows 26/35;
%   faces 35/50; control next_to 37/44 = 84%; near misses 35/53 = 66%; clear negatives 36/40 = 90%.
% * GPT-6 98 / 96 / 90 (n = 536 / 248 / 349): CLAIMS §4 "E5-fail", hard/e5_fail/REPORT.md.  T2 = labels + positions (oracle perception).
% * Surface-gap correction, 9.5% of closest / farthest claims: RESULTS §4.5 (11-12).
% * Writer's own relations vs the reader's, 15--25 points: RESULTS §1.6 (generation mode); quoted in sec/5.
% Parked for the supplement: q-sensitivity (3.5 / 7.2 / 14.5% at q = 0.05 / 0.10 / 0.20), text-channel loss by relation family,
% unbound-reason distribution, verifier revision history, per-family human table.
\section{Scoring exactness}
\label{sec:scorer}

\paragraph{Definition.}
Let $T$ be the target set of a sentence, $C \supseteq T$ its candidates and $D = C \setminus T$ its distractors (\cref{sec:pipeline}).
The reader turns the sentence into predicates $p_1,\dots,p_K$ over object ids; the verifier evaluates each on the scene and returns $\hat{E}_k \subseteq C$, the candidates that satisfy it.
The sentence is \emph{certified} when
\begin{equation}
  K \ge 1 \quad\text{and}\quad \bigcap_{k=1}^{K} \hat{E}_k = T .
  \label{eq:exact}
\end{equation}
It can fail in two ways: a distractor survives every predicate, or one predicate excludes a target.
A sentence that coordinates groups (\emph{the two by the window and the one at the door}) is read in parts; predicates conjoin within a part and the parts are united (\supp{supp:parts}).
The condition $K \ge 1$ is deliberate.
A sentence that offers nothing to check is not certified, even when its target is the only object of its class and the class name would do: we score what the sentence says about space.

\paragraph{The reader is blind.}
The reader is GPT-5.5 with the sentence and the room's instance list.
It does not see the target, the candidates, our evidence or the relations the writer reported, and for every clause it cannot bind to object ids it must give the reason: no viewpoint, ambiguous anchor, object not in the list, not spatial.
We keep it blind because a reader that knows the answer reads toward it.
Given a \emph{wrong} target on purpose, it certified that object in 2 of 15 trials, in one of them by silently dropping the two clauses that excluded it.
Blindness costs little: in an earlier run of 200 pairs, telling the reader the true target raised the rate of the original sentences by one point.

\paragraph{The verifier contains no model.}
It evaluates 42 predicates on the instance table and the points: proximity (\emph{next to}, \emph{closest to}, \emph{touching}, \emph{between}), support and vertical order, places in the room (wall, corner, centre), directions from a stated standpoint, rows and positions in them, and the faces and ends of large furniture (\supp{supp:vocab}).
The same input always gives the same verdict.
With the operators that produce the writer's evidence it shares only a few measuring primitives---surface gap, support, collinearity, row decomposition; above these, neither calls the other.
An error in a primitive therefore shows on both sides, while an error above them stays on one, so the verifier's output is audited on its own, below.

\paragraph{Strict and lenient.}
A predicate the reader extracts can be false of the target itself.
The \emph{strict} score intersects it like any other, and the sentence fails; the \emph{lenient} score skips it.
The two differ in this rule and nothing else.
A skipped predicate may be the writer's mistake, the reader's over-translation (\emph{along the wall} read as \emph{against the wall}) or the verifier's rigidity (a wall test in a room that is not a rectangle).
We cannot yet separate the three, so we report both and take the strict score as the measurement.

\paragraph{A lower bound.}
Write $E_k$ for the candidates that truly satisfy $p_k$.

\noindent\textbf{Proposition 1.} \emph{Suppose that for every $k$ the verifier has no false positive on a target, $\hat{E}_k \cap T \subseteq E_k$, and no false negative on a distractor, $E_k \cap D \subseteq \hat{E}_k$. Then $\bigcap_k \hat{E}_k = T$ implies $\bigcap_k E_k = T$.}

\noindent Under these assumptions every strictly certified sentence is truly exact, and the certified rate bounds the true rate from below (proof and sensitivity analysis in \supp{supp:proof}).
The assumptions concern opposite errors on complementary sets, so no overall accuracy of the verifier implies either of them.
They will not hold perfectly, and the proposition takes the predicates as given; it says nothing about what the reader failed to recover.
We therefore also measure the scorer against people.

\paragraph{Against people.}
Annotators saw a scene as an interactive point cloud with the candidates lettered at random, and one sentence, without knowing whether it was an original or a rewrite; they selected every object the sentence fits (\cref{tab:human}).
Three things follow.
The scorer is conservative: people resolve \TODO{47}\% of the originals and \TODO{70}\% of the rewrites, where the scorer certifies \TODO{12}\% and \TODO{51}\% of the same items.
A sentence it does not certify is one whose uniqueness cannot be established from its spatial relations; a person may still find the object by its colour or by common sense, which the verifier does not see.
A certificate is worth about a second opinion: annotators resolve \TODO{75}\% of certified sentences, and when one annotator resolves a sentence a second does so \TODO{79}\% of the time.
And people and scorer order the sentences alike: human success falls as same-class objects multiply, and falls furthest for the originals.

We also checked the verifier predicate by predicate, on the structures we define ourselves: rows, furniture faces and stacks.
Annotators judged single statements with the objects marked, half of them made false by pointing at an object the verifier excludes.
They agree with the verifier on \TODO{78}\% of the statements they could decide (\TODO{84}\% for plain adjacency, the control), and disagreement concentrates on near misses: \TODO{66}\% agreement there against \TODO{90}\% on clear cases.
The verifier has been wrong before in exactly this way.
Until a human reader caught it, it measured gaps between boxes instead of surfaces, which changed the winner in 9.5\% of \emph{closest} and \emph{farthest} claims; the certified rates in \cref{sec:data,sec:model} all follow the correction (\supp{supp:ruler}).

A reading by a different machine agrees in order.
On distractor-heavy targets, GPT-6 reads the rewrite with every object's label and position at hand and nothing from the scorer; it resolves 98\% of the strictly certified sentences, 96\% of those certified only leniently and 90\% of the rest (\cref{sec:model}).
A rejection means \emph{not certifiable from the sentence}, not \emph{unsolvable}.
% Source: /data/sujie/horus/_rerun/e7 (task A1), computed by Seshat-Star/analysis/_e7_snapshot.py.
% SNAPSHOT 11-16 21:33 CST: 536 judgements, 4 annotators, 380 of 400 (item, arm) units; annotation still running.
%   all:           orig 127/271 = 46.9   rewrite 185/265 = 69.8
%   certified:     orig  22/32  = 68.8   rewrite 107/141 = 75.9   (both arms 129/173 = 74.6)
%   not certified: orig 105/239 = 43.9   rewrite  78/124 = 62.9   (both arms 183/363 = 50.4)
%   same-class objects <=2: 52/72, 64/68;  3-5: 43/70, 52/67;  6-9: 23/65, 38/69;  >=10: 9/64, 31/61
% "cannot tell" / "none fits" count as a miss (user, 11-16).  Re-run the script and refresh when annotation closes.
\begin{table}[t]
\centering
\small
\caption{\textbf{People against the scorer.} Share of judgements (\%) in which an annotator selects exactly the target set, on 200 validation items shown as the original or as the rewrite; ``cannot tell'' counts as a miss. \TODO{Snapshot: 536 judgements by four annotators; refresh when annotation closes.}}
\label{tab:human}
\begin{tabular}{@{}lcc@{}}
\toprule
Annotator resolves the sentence & Original & Rewrite \\
\midrule
All items & 47 & 70 \\
\addlinespace
Scorer certifies it (strict score) & 69 & 76 \\
Scorer does not & 44 & 63 \\
\addlinespace
Same-class objects: at most 2 & 72 & 94 \\
\quad 3--5 & 61 & 78 \\
\quad 6--9 & 35 & 55 \\
\quad 10 or more & 14 & 51 \\
\bottomrule
\end{tabular}
\end{table}

